%--------------------------------------------------------------------
\newpage
\section{数据传递}
\label{sec:transfers}

\tkzNamePack{tkz-elements} 宏包严格地把\emph{数学计算}与\emph{图形渲染}分离开来。所有几何对象都在 Lua 内部创建和操作，而绘图则由 \tkzNamePack{TikZ} 或 \tkzNamePack{tkz-euclide} 完成。

所谓\emph{传递}（transfer），就是把数据从 Lua 计算层显式地传送到 \TeX\ 渲染层的操作。

\medskip
从概念上讲，其架构可概括为：

\begin{center}
计算（Lua）$\longrightarrow$ 传递 $\longrightarrow$ 渲染（TikZ/\TeX）
\end{center}

在被显式传递之前，对象对 \TeX\ 是不可见的。这种设计保证了数值稳健性，并在计算与呈现之间保持清晰的分离。

%--------------------------------------------------------------------
\subsection{概念概览}

\tkzNamePack{tkz-elements} 提供三种传递机制：

\begin{enumerate}
\item 通过 \TeX\ 输入流的即时传递；
\item 通过专用宏的基于对象的传递；
\item 通过外部数据文件的基于文件的传递。
\end{enumerate}

每种机制适用于不同的场合，取决于所传递数据的规模和性质。

\subsection{通过 \texttt{tostring} 进行对象序列化}
\label{subsec:transfers_tostring}

宏 \tkzNameMacro{tkzUseLua} 是从 Lua 到 \TeX\ 的基本桥梁。它对一个 Lua 表达式求值，并把结果打印到 \TeX\ 输入流中。在内部，它使用 Lua 的 \texttt{tostring()} 机制：

\begin{verbatim}
\def\tkzUseLua#1{\directlua{tex.sprint(#1)}}
\end{verbatim}

因此，任何 \tkzNamePack{tkz-elements} 对象都可以定义一个 Lua \texttt{\_\_tostring} 元方法来控制它在 \TeX\ 中的文本表示。这一点尤其用于复数（\tkzClass{point}）。

\medskip
\texttt{点的示例：}
\begin{tkzexample}[latex=.5\textwidth]
\directlua{
  init_elements()
  z.A = point(1, -2)
}
\tkzUseLua{tostring(z.A)}
\end{tkzexample}

于是，与点 $A$ 关联的\emph{复数附标}就通过 Lua 的 \texttt{\_\_tostring} 元方法（经由 \texttt{tostring}）打印出来。

\medskip
\texttt{矩阵的示例： }

宏 \tkzNameMacro{tkzUseLua} 把一个 Lua 表达式的\emph{值}作为文本打印出来（内部经由 \texttt{tostring(<expr>)}）。%
\footnote{因此，对象可以通过定义一个 Lua \texttt{\_\_tostring} 元方法来自定义它在行文中的表示。}


\medskip
\begin{tkzexample}[latex=.5\textwidth]
\directlua{
  init_elements()
  local a, b, c, d = 1, 2, 3, 4
  M.new = matrix({ { a, b }, { c, d } })}
  \tkzUseLua{M.new:print()}
\end{tkzexample}

\texttt{其他情况}
要传递 \code{nil}、\code{true} 和 \code{false}，必须使用 \code{tostring}；不过，数字或字符串可以直接传递。

\begin{verbatim}
  \tkzUseLua{tostring(nil)}
  \tkzUseLua{tostring(true)}
  \tkzUseLua{"TANGENT"}
  \tkzUseLua{math.pi}
\end{verbatim}
%--------------------------------------------------------------------
\subsection{即时传递}
\label{subsec:transfer_immediate}

当 Lua 使用 \texttt{tex.print} 直接把数据打印到 \TeX\ 输入流时，就发生即时传递。

\subsubsection*{数值示例}

\begin{tkzexample}[latex=.5\textwidth]
\directlua{
  init_elements()
  z.A = point(3,4)
  tex.print(point.abs(z.A))
}
\end{tkzexample}

计算所得的值被直接插入文档。

\subsubsection*{布尔示例}

\begin{tkzexample}[latex=.5\textwidth]

\directlua{
  init_elements()
  z.A = point(0, 0)
  z.B = point(3, 0)
  z.C = point(0, 2)
  T.ABC = triangle(z.A, z.B, z.C)
  tex.print(tostring(T.ABC:check_equilateral()))
}
\end{tkzexample}

这一机制让 Lua 计算能够通过条件语句影响 \TeX\ 的逻辑。

\medskip
\texttt{注意： } 可以自行决定何时进行传递。为此专门创建了一个宏，但你也可以自建：它叫做 \tkzcname{tkzUseLua}。其定义如下：

\begin{mybox}
  \begin{verbatim}
  \def\tkzUseLua#1{\directlua{tex.print(#1)}}
\end{verbatim}
\end{mybox}

该宏把一个 Lua 变量或表达式的值直接打印到 TeX 流中。

\medskip
\texttt{示例。}

下面的 Lua 代码计算两条直线是否相交：

\directlua{
 z.a = point(4, 2)
 z.b = point(1, 1)
 z.c = point(2, 2)
 z.d = point(5, 1)
 L.ab = line(z.a, z.b)
 L.cd = line(z.c, z.d)
 det = (z.b - z.a) ^ (z.d - z.c)
 if det == 0 then bool = true
   else bool = false
 end
 x = intersection (L.ab, L.cd)}

两条直线的交点位于附标为：\tkzUseLua{x} 的点处。

\vspace{1em}
\begin{minipage}{0.5\textwidth}
\begin{verbatim}
\directlua{
  z.a = point(4, 2)
  z.b = point(1, 1)
  z.c = point(2, 2)
  z.d = point(5, 1)
  L.ab = line(z.a, z.b)
  L.cd = line(z.c, z.d)
  det = (z.b - z.a) ^ (z.d - z.c)
  if det == 0 then bool = true
    else bool = false
  end
  x = intersection (L.ab, L.cd)}

两条直线的交点位于附标为：
\tkzUseLua{tostring(x)} 的点处。


\begin{tikzpicture}
  \tkzGetNodes
  \tkzInit[xmin =0,ymin=0,xmax=5,ymax=3]
  \tkzGrid\tkzAxeX\tkzAxeY
  \tkzDrawPoints(a,...,d)
  \ifthenelse{\equal{\tkzUseLua{bool}}{%
  true}}{\tkzDrawSegments[red](a,b c,d)}{%
  \tkzDrawSegments[blue](a,b c,d)}
  \tkzLabelPoints(a,...,d)
\end{tikzpicture}
\end{verbatim}
\end{minipage}
\begin{minipage}{0.5\textwidth}
\begin{center}
\begin{tikzpicture}
  \tkzGetNodes
  \tkzInit[xmin =0,ymin=0,xmax=5,ymax=3]
  \tkzGrid\tkzAxeX\tkzAxeY
  \tkzDrawPoints(a,...,d)
  \ifthenelse{\equal{\tkzUseLua{bool}}{true}}{
  \tkzDrawSegments[red](a,b c,d)}{%
  \tkzDrawSegments[blue](a,b c,d)}
  \tkzLabelPoints(a,...,d)
\end{tikzpicture}
\end{center}
  \end{minipage}

%----------------------------------------------
% ---> 在此处插入额外的自定义示例
%----------------------------------------------

%--------------------------------------------------------------------
\subsection{基于对象的传递}
\label{subsec:transfer_objects}

存放在 Lua 表中的几何对象，必须先转换成与 TikZ 兼容的结构，才能绘制。

\subsubsection{传递点}
\label{transferring_points}

点在内部以复数附标的形式存放。宏 \tkzNameMacro{tkzGetNodes} 把它们转换成 TikZ 坐标。

\begin{tkzexample}[latex=7cm]
  \directlua{
     init_elements()
     z.o   = point(0, 0)
     z.a_1 = point(2, 1)
     z.a_2 = point(1, 2)
     z.ap  = z.a_1 + z.a_2
     z.app = z.a_1 - z.a_2
  }

  \begin{center}
    \begin{tikzpicture}[ scale = 1.5]
       \tkzGetNodes
       \tkzDrawSegments(o,a_1 o,a_2 o,a' o,a'')
       \tkzDrawSegments[red](a_1,a' a_2,a')
       \tkzDrawSegments[blue](a_1,a'' a_2,a'')
       \tkzDrawPoints(a_1,a_2,a',o,a'')
       \tkzLabelPoints(o,a_1,a_2,a',a'')
    \end{tikzpicture}
  \end{center}
\end{tkzexample}

\subsubsection{传递路径与曲线}


\begin{tkzexample}[latex=.4\textwidth]
\directlua{
 init_elements()
 PF.ex = pfct("(cos(t))^3", "(sin(t))^3")
 PA.curve = PF.ex:path(0, 2*math.pi, 100)
 }

\begin{center}
 \begin{tikzpicture}[scale=2]
  \tkzDrawCoordinates[smooth,red,fill=orange](PA.curve)
\end{tikzpicture}
 \end{center}
\end{tkzexample}

该 Lua 对象被序列化为一个由 TikZ 解释的坐标列表。

%----------------------------------------------
% ---> 在此处插入高级 fct / pfct 示例
%----------------------------------------------

%--------------------------------------------------------------------
\subsubsection{基于文件的传递}


对于大型数据集或高分辨率曲线，生成一个外部文件可能更为合适。

这一次，传递将借助一个外部文件来完成。下面这个示例基于上一个示例，但改用一张表。

\texttt{生成数据文件： } 文件 \texttt{tmp.table} 由函数 $f$ 创建。

文件 \texttt{tmp.table} 现在包含数值数据。
 \texttt{使用生成的文件： } 这里使用 \tkzNamePack{TikZ} 绘制曲线。




\begin{tkzexample}[latex=.5\textwidth]
 \directlua{
  init_elements()
   z.a = point(1, 0)
   z.b = point(3, 2)
   z.c = point(0, 2)
   A,B,C = tkz.parabola (z.a, z.b, z.c)

 function f(t0, t1, n)
  local out=assert(io.open("tmp.table","w"))
  local y
  for t = t0,t1,(t1-t0)/n  do
   y = A*t^2+B*t +C
   out:write(utils.checknumber(t), " ",
             utils.checknumber(y), " i\string\n")
  end
  out:close()
 end}

\begin{tikzpicture}
   \tkzGetNodes
   \tkzInit[xmin=-1,xmax=5,ymin=0,ymax=5]
   \tkzDrawX\tkzDrawY
   \tkzDrawPoints[red,size=2](a,b,c)
   \directlua{f(-1,3,100)}%
   \draw[domain=-1:3] plot[smooth]
        file {tmp.table};
\end{tikzpicture}
 \end{tkzexample}




基于文件的传递提供了可扩展性，并能与 \texttt{pgfplots} 等外部绘图工具互通。

%----------------------------------------------
% ---> 在此处插入高分辨率示例
%----------------------------------------------

%--------------------------------------------------------------------
\subsection{\TeX--Lua 的动态交互}
\label{subsec:dynamic_transfer}

传递是双向的。
\TeX\ 可以把参数传给 Lua 进行计算。

\medskip
\begin{verbatim}
\newcommand{\ComputeSquare}[1]{%
  \directlua{tex.print(#1 * #1)}%
}
\end{verbatim}

\medskip
用法：
\newcommand{\ComputeSquare}[1]{%
  \directlua{tex.print(#1 * #1)}%
}
\begin{verbatim}
5 的平方是 \ComputeSquare{5}。
\end{verbatim}

\medskip
5 的平方是 \ComputeSquare{5}。

\medskip

这种交互实现了动态的参数化几何，文档输入可直接控制 Lua 计算。

%----------------------------------------------
% ---> 在此处插入动态几何示例
%----------------------------------------------

\medskip
因此，\tkzNamePack{tkz-elements} 中的传递可以是即时的、基于对象的或基于文件的，取决于数据的规模和结构。这种分层架构保证了计算与图形表示之间的清晰分离。
